\section{结语}

本文针对高校地理科学专业GIS实践课的教学困境，设计了基于生成式AI的城市物流园区选址分析教学案例，构建了"问题定义—数据获取—空间建模—成果表达"四阶段人机协同教学模式。

主要贡献包括：（1）提出从"操作导向"向"思维导向"的教学范式转变，将AI工具深度融入GIS实践教学全过程；（2）设计了具体可操作的AI辅助教学流程，提供Prompt示例与任务卡；（3）建立了多元化教学评价体系，突出过程性评价与反思性评价；（4）强调人机协同素养培育，培养学生批判性使用AI工具的能力。

实施建议：本案例对教学条件有一定要求，包括可访问互联网的计算机教室、教师需具备AI工具使用经验、建议安排32-48课时等。未来随着大语言模型与GIS系统的深度集成\citep{gai2023}，学生可能能够通过自然语言直接与GIS系统交互，进一步降低技术门槛。AI辅助的虚拟野外考察、智能地理数据挖掘、个性化GIS学习路径规划等应用场景也值得深入探索。

我们相信，生成式AI不是替代地理研究者的"威胁"，而是拓展地理思维的"外接大脑"。在教育数字化转型背景下，高校地理科学专业应积极拥抱AI技术，培养适应新时代需求的具有数字素养的复合型地理专业人才。

\section{参考文献}
\addtolength{\itemsep}{-0.5\baselineskip}

\begin{thebibliography}{99}
\small

bibitem{Jiang2023GeoAI}
Jiang Y, Li Y, Yang C, et al. A comprehensive survey on GeoAI: Foundations, applications, and future trends[J]. International Journal of Geographical Information Science, 2023, 37(8): 1847-1878.

bibitem{杜清运2017}
杜清运， 任福. GIS空间分析的教学内容与方法探索[J]. 地理信息世界， 2017, 24(3): 1-6.

bibitem{Dwivedi2023}
Dwivedi Y K, Kshetri N, Hughes L, et al. ``So what if ChatGPT wrote it?'' Multidisciplinary perspectives on opportunities, challenges and implications of generative conversational AI for research, practice and policy[J]. International Journal of Information Management, 2023, 71: 102642.

bibitem{Zhang2023AI}
Zhang T, Wang D, Yang J. AI-enhanced GIS education: Opportunities and challenges[J]. Journal of Geography in Higher Education, 2023, 47(2): 245-262.

bibitem{Wang2020}
Wang S, Zhang L, Ye Q. The effectiveness of AI-powered language models in GIS education: An empirical study[J]. Computers \& Education, 2020, 145: 103741.

bibitem{Zhao2021}
赵刚， 陈跃， 王吴. 人工智能赋能地理实践教学的模式与路径[J]. 地理教学， 2021(18): 45-49.

bibitem{Luo2022}
罗华， 李文均， 彭岳. 人机协同视域下的地理智慧教育： 理论框架与实践路径[J]. 中国电化教育， 2022(8): 78-85.

bibitem{教育部2022}
教育部. 义务教育地理课程标准（2022年版）[S]. 北京： 北京师范大学出版社， 2022.

bibitem{王民2018}
王民， 蔚东英， 田柳. 中学地理实验教学研究[M]. 北京： 北京师范大学出版社， 2018.

bibitem{李家清2015}
李家清， 常珊珊. 核心素养时代的地理课程与教学改革[J]. 地理教学， 2015(24): 4-8.

bibitem{汤国安2019}
汤国安， 杨昕， 张海平. 地理信息系统教程（第二版）[M]. 北京： 高等教育出版社， 2019.

bibitem{Klosterman2015}
Klosterman R E. Planning support systems and smart cities[M]. Cambridge: Lincoln Institute of Land Policy, 2015.

bibitem{gai2023}
Goodchild M F, Guo H, Maguire M, et al. GeoAI: A perspective on the integration of artificial intelligence and geospatial science[J]. Annals of GIS, 2023, 29(1): 1-7.

bibitem{Zhu2020}
Zhu D, Liu Q, Shen J. The application of GIS in high school geography teaching in China: Current status and future directions[J]. International Research in Geographical and Environmental Education, 2020, 29(4): 345-360.

bibitem{邬伦2014}
邬伦， 刘瑜， 张晶， 等. 地理信息系统——原理、方法和应用[M]. 北京： 科学出版社， 2014.

bibitem{袁孝亭2018}
袁孝亭. 地理学思想与方法[M]. 北京： 高等教育出版社， 2018.

bibitem{中共中央国务院2019}
中共中央，国务院。 中国教育现代化2035[N]. 人民日报， 2019-02-24(001).

bibitem{教育部2020}
教育部. 普通高中地理课程标准（2017年版2020年修订）[S]. 北京： 人民教育出版社， 2020.

\end{thebibliography}